% \endlinechar governs what a line end contributes -- including inside a macro
% argument, which is where it was being ignored: argument reading used a
% hardcoded 13 whatever the document had set.
%
% Only forms verified in isolation against pdftex are here. Cases that change
% \endlinechar and then span lines are genuinely delicate -- a line end that
% contributes nothing merges the next line into this one -- and two attempts at
% writing them produced files that tested something other than what I intended.
\input plain
\catcode`\@=11
\def\arg#1{(#1)}
\message{<default=\arg{a
b}>}
\endlinechar=`\!
\message{<bang-arg=\arg{a
b}>}
\message{<bang-plain=x
y>}
\endlinechar=13
% expl3's regime: the line end is a space *character*, and the space character
% is ignored -- so line ends vanish completely, even inside arguments.
\catcode`\ =9 \endlinechar=32
\def\outer@arg#1{[#1]}
\message{<expl-arg=\outer@arg{a
b}>}
% A \def read as part of an argument keeps its parameter text intact: no
% phantom space delimiter from the line end. This is what stopped expl3.
\outer@arg{\def\inner#1
  {(#1)}}
\message{<expl-def=\inner X>}
\endlinechar=13 \catcode`\ =10
\bye
